\Needspace{25\baselineskip}
\section{The lower-mass quotient}
\label{sec:lower-mass-quotient}

From now on $X$ is again an arbitrary connected compact complex manifold of dimension $n\geq1$, and $\omega$ is a Hermitian form on $X$. By \cref{thm:gaud-volume-mass-extrema}, the lower mass $v_-(\omega)$ is the infimum of the volumes $\mathcal V_\mu(\omega)$ over the set $\mathscr M$ of smooth strictly positive probability measures, and \eqref{eq:gaud-lower-mass-formula} expresses it through Gauduchon metrics alone. We relate attainment of this infimum to smooth minimizers of the mass, and then extend the quotient to nef backgrounds.

\begin{proposition}\label{prop:lower-mass-smooth-attainment}\index{Monge--Ampere mass@Monge--Amp\`ere mass!smooth minimum}
The infimum in \eqref{eq:gaud-lower-mass-formula} is attained if and only if the infimum of the masses $\int_X(\omega+\ddc v)^n$ over smooth functions $v$ with $\omega+\ddc v>0$ is attained. For every such minimizing potential $v$, the metric $\omega+\ddc v$ is Gauduchon.
\end{proposition}
\begin{proof}
Suppose that a Gauduchon metric $\gamma$ attains the infimum in \eqref{eq:gaud-lower-mass-formula}, and put $m=v_-(\omega)^{1/n}$. For the measure $\mu_\gamma=\gamma^n/\int_X\gamma^n$, equality holds in the H\"older inequality of the proof of \cref{thm:gaud-volume-mass-extrema}, so the test $\gamma$ gives
$\lambda_{\mu_\gamma}(\omega)\leq m$.
\Cref{thm:gaud-volume-mass-extrema} gives the opposite inequality. Hence the solution of the smooth equation for $\mu_\gamma$ has mass $m^n=v_-(\omega)$, which gives a minimizing smooth potential. For $n\geq2$, $\gamma$ also minimizes the quotient for $\mu_\gamma$, so by \cref{thm:gaud-volume-dual-formula} the solution form is conformal to $\gamma$. Since both have volume forms proportional to $\mu_\gamma$, the conformal factor is constant, and the solution metric is Gauduchon.

Conversely, let $v$ be a smooth potential minimizing the mass. For every smooth real $f$, the potentials $v\pm tf$ are admissible for sufficiently small $t>0$. The vanishing of the first variation of the mass therefore gives
\[
 \int_X f\,\ddc\bigl((\omega+\ddc v)^{n-1}\bigr)=0
 \quad\hbox{for every smooth real }f.
\]
Thus $\Omega=\omega+\ddc v$ is Gauduchon. Taking $\gamma=\Omega$ in the quotient and using Stokes' theorem for the Gauduchon metric $\Omega$ gives
\[
 \frac{\int_X\omega\wedge\Omega^{n-1}}
      {(\int_X\Omega^n)^{(n-1)/n}}
             =(\int_X\Omega^n)^{1/n}=v_-(\omega)^{1/n},
\]
so the infimum in \eqref{eq:gaud-lower-mass-formula} is attained. In dimension one all metrics are Gauduchon and all masses equal $\int_X\omega$.
\end{proof}

\Needspace{9\baselineskip}
The infimum over densities also makes sense beyond the positive cone. For a pseudoeffective smooth real $(1,1)$-form $\theta$ on $X$, put
\[
 \lambda_*(\theta)=
 \inf_{\gamma\text{ Gauduchon}}
 \frac{\int_X\theta\wedge\gamma^{n-1}}
      {(\int_X\gamma^n)^{(n-1)/n}}.
\]
The H\"older argument in the proof of \cref{thm:gaud-volume-mass-extrema} proves
\[
 \lambda_*(\theta)=\inf_{\mu\in\mathscr M}\lambda_\mu(\theta).
\]
Like each $\lambda_\mu$, the function $\lambda_*$ is finite, nonnegative, concave, positively homogeneous, and invariant under smooth $\ddc$ changes of $\theta$. Suppose $\theta$ is nef, fix a Hermitian form $h$, and for real $\epsilon>0$ choose a Hermitian representative $\omega_\epsilon$ of $\theta+\epsilon h$. Interchanging the infima over the test metric and over $\epsilon>0$, as in \cref{thm:gaud-volume-nef-limit}, gives
\[
 \lambda_*(\theta)
 =\lim_{\epsilon\downarrow0}\lambda_*(\theta+\epsilon h)
 =\lim_{\epsilon\downarrow0}v_-(\omega_\epsilon)^{1/n}.
\]
By invariance under smooth $\ddc$ changes, the last expression does not depend on the choice of the representatives $\omega_\epsilon$. This extends the lower-mass formula to the nef boundary. We do not compare $\lambda_*$ with other notions of lower volume for singular classes.
